\documentclass{article}
\usepackage[portuguese]{babel}
\usepackage{listings}
\usepackage{xcolor}

\definecolor{codegreen}{rgb}{0,0.6,0}
\definecolor{codegray}{rgb}{0.5,0.5,0.5}
\definecolor{codepurple}{rgb}{0.8,0,0.2}
\definecolor{backcolour}{rgb}{.8,.8,1}

\lstdefinestyle{mystyle}{
    backgroundcolor=\color{backcolour},   
    commentstyle=\color{codegreen},
    keywordstyle=\color{blue},
    numberstyle=\tiny\color{codegray},
    stringstyle=\color{codepurple},
    basicstyle=\ttfamily\footnotesize,
    breakatwhitespace=false,         
    breaklines=true,                 
    captionpos=b,                    
    keepspaces=true,                 
    numbers=left,                    
    numbersep=5pt,                  
    showspaces=false,                
    showstringspaces=false,
    showtabs=false,                  
    tabsize=2
}

\lstset{style=mystyle}
\begin{document}

\section{AA7}

Escreva um programa, em linguagem Python, para ler n notas de uma disciplina, via console, armazená-los num vetor e fazer as seguintes tarefas:
\begin{enumerate}
    \item imprimir o vetor na ordem lida;
    \item imprimir o vetor em ordem crescente;
    \item imprimir a maior e a menor nota;
    \item imprimir a nota média;
    \item calcular e imprimir o desvio padrão;
    \item imprimir o número de alunos que estão abaixo de um desvio padrão abaixo da média;
    \item plotar um gráfico de dispersão unidimensional para os dados.
    
\end{enumerate}

O programa inicia obtendo o numero n de notas do usuario.Logo apos as notas sao obtidas do usuario e armazenadas no vetor(lista). 
\begin{lstlisting}[language=python]
n= int(input("Digite o numero de notas:\n"))
notas = []
for i in range(n):
    notas.append(int(input(f"Digite a {i+1}a nota")))
print(f"Lista de notas: {notas}")
\end{lstlisting}

\subsection{Ordenar o vetor}
Essa e a parte mais dificil. Para ordenar esse vetor vamos encontrar o elemento de menor valor e colocar na primeira posicao. Na proxima itracao, ignoramos a primeira posicao e procuramos o menor valor dentre os que sobraram e o colocamos proximo ao primeiro e assim por diante. 

\begin{lstlisting}[language=Python]
 ordenado = list(notas) #Criamos uma copia para
    for i in range(n):  #nao alterar o vetor inicial
        posicaoMenor = i
        for j in range(i+1,n):
            if (ordenado[j] < ordenado[posicaoMenor]):
                posicaoMenor = j
        aux = ordenado[i]
        ordenado[i] = ordenado[posicaoMenor]
        ordenado[posicaoMenor] = aux
    print(f"Lista de notas ordenada: {ordenado}")
\end{lstlisting}

\subsection{Maior e menor nota}
Como o vetor ja foi ordenado a maior nota sera o ultimo item da lista, e o menor sera o primeiro item, Assim nao precisaremos iterar sobre a lista para procurar.Dessa forma iremos apenas selecionar e imprimir.

\begin{lstlisting}[language=Python]
 maior = ordenado[-1]
menor = ordenado[0]
 print(f"Maior nota: {maior}, menor: {menor}")
\end{lstlisting}

\subsection{Media, desvio padrao e numero de notas baixas}
O algoritmos para definir a media e o desvio padrao sao semelhantes. Um laco for precorre o vetor acumula valores em uma variavel.
\[
M_a = \frac{\sum x_i}{n}, \sigma = \frac{\sum(x_i-Ma)^2}{n}
\]
O que esta dentro do somatorio e o mesmo que esta dentro do laco for.
\begin{lstlisting}[language=Python]
media = 0
for nota in notas:
    media+=nota
media = media/n
print(f"Media: {media}")

desvio = 0
for nota in notas:
    desvio += (nota - media)**2
desvio = desvio/n
desvio = desvio**.5
print(f"Desvio padrao: {desvio}")   
\end{lstlisting}
Para obter o numero de notas baixas, percorremos o vetor de notas e toda vez que a nota for menor que $M_a - \sigma$ adicionamos 1 a variavel correspondente.

\begin{lstlisting}[language=Python]
notasBaixas = 0
for nota in notas:
    if(nota <= media-desvio):
        notasBaixas+=1
print(f"Notas menores que a media - desvio padrao:{notasBaixas}\n")
\end{lstlisting}

\subsection{Plotar grafico com valor das notas}
Aqui a biblioteca faz quase tudo. precisamos apenas instalar a biblioteca, importar no inicio do codigo, fornecer os dados e pronto.

\begin{lstlisting}[language=Python]
import matpllotlib.pyplot as mp
#[...
#...]
mp.hist(notas)
mp.show()
\end{lstlisting}
\subsection{Codigo completo:}
\begin{lstlisting}[language=Python]
import matplotlib.pyplot as mp

print("AA-07-Marcus Sousa\n\n")
deveParar = ""

while (deveParar == ""):
    n= int(input("Digite o numero de notas:\n"))
    notas = []
    for i in range(n):
        notas.append(int(input(f"Digite a {i+1}a nota")))
    print("\n-------00-------\n")
    print(f"Lista de notas: {notas}")

    ordenado = list(notas)
    for i in range(n):
        posicaoMenor = i
        for j in range(i+1,n):
            if (ordenado[j] < ordenado[posicaoMenor]):
                posicaoMenor = j
        aux = ordenado[i]
        ordenado[i] = ordenado[posicaoMenor]
        ordenado[posicaoMenor] = aux
    print(f"Lista de notas ordenada: {ordenado}")

    maior = ordenado[-1]
    menor = ordenado[0]

    print(f"Maior nota: {maior}, menor: {menor}")

    media = 0
    for nota in notas:
        media+=nota
    media = media/n
    print(f"Media: {media}")

    desvio = 0
    for nota in notas:
        desvio += (nota - media)**2
    desvio = desvio/n
    desvio = desvio**.5
    print(f"Desvio padrao: {desvio}")   

    notasBaixas = 0
    for nota in notas:
        if(nota <= media-desvio):
            notasBaixas+=1
    print(f"Notas menores que a media - desvio padrao:{notasBaixas}\n")
    mp.hist(notas)
    mp.show()

    deveParar = input("Deseja continuar?(enter para continuar, qualquer tecla para parar)\n")

print("Ate mais")
\end{lstlisting}
\end{document}
